\begin{proposition}[Euler signs for two of the rational boundary classes]\label{cdp:st:C.2.RAT.1}
\Uses{the hypotheses, notation and definitions stated here; \cref{cdp:st:A.5.16}; \cref{cdp:eq:C.RA.7}; \cref{cdp:eq:C.RA.8}; \cref{cdp:eq:C.RA.9}; \cref{cdp:eq:C.RA.10}; \cref{cdp:st:A.5.5}; \cref{cdp:st:A.5.15}; \cref{cdp:st:A.5.17}; \cref{cdp:eq:C.RA.14}; \cref{cdp:eq:C.RA.15}; \cref{cdp:eq:C.RA.6}; \cref{cdp:eq:C.RA.13}; \cref{cdp:eq:C.RA.2}}


Let $f:Y\to B$ be a proper holomorphic map from a smooth compact
complex threefold with $b_2(Y)=0$, with connected fibres and smooth
general fibre a complex two-torus. Let $mS$ be a full fibre with
irreducible nonnormal reduced support. Write
$h=\eta\pi:Z\to T\to S$ for the minimal resolution of its
normalization and suppose that $Z$ is rational. Use the effective
divisor $K_Z+D=0$ of \ref{cdp:st:A.5.16}.

If a contraction to $\mathbf P^2$ sends $D$ to three concurrent
reduced lines, then $e(S)\ge1$. More precisely,
$e(S)=-1+N_0+e(\Gamma_{\rm red})\ge1+N_0$, with $N_0$ as in
(A.RT.12).

Consider also a specified morphism $\rho:Z\to\mathbb F_n$,
$n\ge1$, which is a sequence of point blowups in reverse and has
\[
 \rho_*D=2C_0+\sum_{j=1}^r a_jF_j,\qquad
 C_0^2=-n,\qquad
 a_j\ge1,\quad \sum_j a_j=n+2,                  \tag{C.RA.1}\label{cdp:eq:C.RA.1}
\]
where the $F_j$ are distinct ruling fibres. Let $C$ be the
strict transform of $C_0$. Suppose $h(C)$ is a curve and the
generic transverse singularity there is a cusp $A_2$.
Then $r\le3$ and $e(S)\ge1$. If $h(C)$ is smooth, then
$r\le2$. These assertions allow singular normalization,
arbitrary vertical conductor multiplicities and a nonreduced
conductor-image scheme.
\end{proposition}
\begin{proof}
\noindent\textbf{Step 1.} Put $G=\Gamma_{\rm red}$. \ref{cdp:st:A.5.16} has already proved
\[
 e(S)=e(Z\setminus D_{\rm red})+e(G).             \tag{C.RA.2}\label{cdp:eq:C.RA.2}
\]
In either boundary class under consideration, $D_{\rm red}$
is simply connected, all its components are rational, and
(A.RT.18)--(A.RT.20) give
\[
 \pi_1(S)=\pi_1(G)\in\{1,\mathbb Z\}.
 \tag{C.RA.3}\label{cdp:eq:C.RA.3}
\]
Write $n_G$ for the number of irreducible components of $G$
and $\beta=b_1(G)\le1$. The normalization maps from their
rational normalizations are homeomorphisms except for the
identification of finitely many points. Consequently
\[
                   e(G)=n_G+1-\beta\ge n_G\ge1.
 \tag{C.RA.4}\label{cdp:eq:C.RA.4}
\]
The exact two-branch incidence argument (A.RT.21) strengthens
this to $e(G)\ge2$ in the present boundary classes.
For three concurrent lines, (A.RT.12)--(A.RT.13) therefore
prove the first assertion.

\noindent\textbf{Step 2.} We prove the cusp assertion. Put $H=h(C)_{\rm red}$.
The transverse $A_2$ curve has a single normalized branch.
Thus $C\to H$ is the normalization map of degree one, and
no other strict conductor component maps onto $H$.
In particular $H^\nu=C=\mathbf P^1$. The restriction
$H^1(G,\OO_G)\to H^1(H,\OO_H)$ is surjective, since its
kernel has one-dimensional support and no second cohomology.
\ref{cdp:st:A.5.16} therefore gives
\[
              \delta:=p_a(H)\le2.              \tag{C.RA.5}\label{cdp:eq:C.RA.5}
\]
Since $H^\nu$ is rational, $\delta$ is the sum of its local
delta invariants. Write $R$ for the union of the other irreducible components of
$G$. The exact constant-sheaf sequence
\[
 0\longrightarrow\mathbb Q_G\longrightarrow
 \mathbb Q_H\oplus\mathbb Q_R\longrightarrow
 \mathbb Q_{H\cap R}\longrightarrow0
\]
gives a surjection
$H^1(G;\mathbb Q)\to H^1(H;\mathbb Q)\oplus H^1(R;\mathbb Q)$,
since $H\cap R$ is finite. Hence $b_1(H)\le b_1(G)\le1$.
If $s$ denotes the number of points of $C$ above singular
points of $H$, normalization of the reduced curve gives
\[
 s=\sum_{p\in\operatorname{Sing}H}r_p
   =b_1(H)+\#\operatorname{Sing}H\le1+\delta
 \quad(\delta>0),\qquad s=0\quad(\delta=0).       \tag{C.RA.6}\label{cdp:eq:C.RA.6}
\]

\noindent\textbf{Step 3.} We need the precise low-delta curve analysis, including its
two possible non-lci cases. A reduced curve germ of delta
invariant at most two has at most three branches.
Three branches require delta two, and the ring is exactly
the subring of its normalization consisting of triples
with equal constant value. This follows by equality of
colengths, both equal to two. Its topological branch
identification alone contributes two loops to a globally
irreducible curve, so $b_1(H)\le1$ excludes it here.

\noindent\textbf{Step 4.} For two branches, the exact sequence for their union gives
$\delta=\delta_1+\delta_2+\ell(\OO_{H_1\cap H_2})$.
Two smooth branches therefore give a node, or, at delta
two, the fibre product matching their constant and linear
terms. The latter has the plane tacnode ring
$\mathbb C\{x,y\}/(y^2-x^4)$.
The only remaining two-branch possibility is a cusp and
a transverse smooth branch, whose ring is
$\mathbb C\{t^2,t^3\}\times_{\mathbb C}\mathbb C\{u\}$.
In embedding dimension three its ideal can be written
\[
                  I=(xz,yz,y^2-x^3).            \tag{C.RA.7}\label{cdp:eq:C.RA.7}
\]
For one branch, the delta invariant equals the number of
gaps of its valuation semigroup. With at most two gaps,
the singular possibilities are
$\langle2,3\rangle$, $\langle2,5\rangle$, and
$\langle3,4,5\rangle$. The first two are plane cusps.
For the last, the conductor is $t^3\mathbb C\{t\}$,
and colength two makes the ring
$\mathbb C+t^3\mathbb C\{t\}$ exactly. Its ideal is
\[
 I=(g_1,g_2,g_3)
   =(y^2-xz,\ yz-x^3,\ z^2-x^2y),\qquad
 (x,y,z)=(t^3,t^4,t^5).                         \tag{C.RA.8}\label{cdp:eq:C.RA.8}
\]
For completeness, the asserted plane form at semigroup
$\langle2,5\rangle$ follows by making an order-two
generator equal to $t^2$ by a change of parameter.
Its conductor is then $t^4\mathbb C\{t\}$, and the ring
is $\mathbb C+\mathbb Ct^2+t^4\mathbb C\{t\}$.
These conductor descriptions also follow in the completed
local rings by eliminating successive higher terms; finite
normalization and faithful flatness of completion give the
same conductor in the convergent rings. Thus there are no
unlisted cases. A germ of embedding dimension two is lci
in the smooth threefold: choose one regular transverse
ambient coordinate in its ideal and then its single plane
equation. We now exclude \eqref{cdp:eq:C.RA.7}--\eqref{cdp:eq:C.RA.8} for our cusp
image.

\noindent\textbf{Step 5.} For \eqref{cdp:eq:C.RA.8}, put
\[
 F=x^5+xy^3+z^3-3x^2yz .
\]
The symbolic square is $I^{(2)}=I^2+(F)$, with
\[
 xF=g_2^2-g_1g_3,\quad
 yF=g_2g_3+xg_1^2,\quad
 zF=g_3^2+xg_1g_2.                              \tag{C.RA.9}\label{cdp:eq:C.RA.9}
\]
The saturation computation $I^2:x^\infty=I^2+(F)$,
including the complete exact Groebner certificate, is
given in \path{code/A_low_delta_cusp.py} and its output.
It computes the symbolic square because the curve is lci
away from the origin and $x$ vanishes on it only at the
origin. The same equality in the analytic local ring
follows by flatness of the convergent extension and the
displayed finite saturation.

\noindent\textbf{Step 6.} Let $N$ be the dual of the torsion-free pullback to $C$ of
the ambient conormal module of $H$. At \eqref{cdp:eq:C.RA.8} the
normalized conormal basis and its dual lifts are
\[
\begin{gathered}
 a=dg_1,\qquad b=dg_2+t\,dg_1,\qquad
 dg_3=-2t^2a+tb,\\
 A=-t^{-3}\partial_z,\qquad
 B=(t^{-5}\partial_y+2t^{-4}\partial_z)/3 .
\end{gathered}                                  \tag{C.RA.10}\label{cdp:eq:C.RA.10}
\]
Their wedge is $a\wedge b=t^6\,\iota_{\gamma'}(dx\wedge
dy\wedge dz)$, where $\gamma=(t^3,t^4,t^5)$.
Thus the determinant formula is
$\det N=K_C\otimes h_C^*K_Y^{-1}(6p)$.
Since every ambient line has degree zero on every compact
curve when $b_2(Y)=0$, its degree here is four.

\noindent\textbf{Step 7.} Write the local equation of $S$ as $cF+f_2$, $f_2\in I^2$.
Its normalized quadratic form is the sum of a regular
quadratic form and $c$ times
\[
 q_F=\frac3t a^2-\frac3{t^2}ab+\frac1{t^3}b^2,
 \qquad \det(q_F)=\frac3{4t^4}.                  \tag{C.RA.11}\label{cdp:eq:C.RA.11}
\]
If $c(0)\ne0$, the determinant has an uncancellable
pole of order four: a cross term with the regular form
has pole at most three. The generic cusp requires rank
one, so $c(0)=0$. Equations \eqref{cdp:eq:C.RA.9} then give
$cF\in I^2$, and the quadratic form is regular.

\noindent\textbf{Step 8.} The intrinsic cubic on its radical has pole order at most
ten at $p$. Here is a direct bound for arbitrary coefficients,
not only for the polynomial $F$. In the lifts \eqref{cdp:eq:C.RA.10},
the differentials $dg_i$ are regular, the second derivatives
$D^2g_i(Au+Bv)^2$ have poles at most ten, and a derivative
of a holomorphic coefficient has pole at most five.
The formula for $D^3(b_{ij}g_i g_j)$ therefore has pole
at most ten. Terms containing an undifferentiated $g_i$
vanish on the curve. The exact derivatives are in the cited computation.
The radical cubic is independent of a tangent component
of the lift: differentiating the identically zero quadratic
form on its radical along the smooth part of the curve
gives this invariance, and equality extends meromorphically.

\noindent\textbf{Step 9.} The same argument is shorter for \eqref{cdp:eq:C.RA.7}. Put
$g_1=xz$, $g_2=yz$, $g_3=y^2-x^3$. Then
$I^{(2)}=I^2+(z g_3)$ and
\[
 x(zg_3)=g_1g_3,\quad y(zg_3)=g_2g_3,\quad
 z(zg_3)=g_2^2-xg_1^2.
 \tag{C.RA.12}\label{cdp:eq:C.RA.12}
\]
On the cusp branch $(t^2,t^3,0)$ use $a=dg_1$, $b=dg_3$.
The quadratic form of $z g_3$ is $t^{-2}ab$, whose
determinant has pole four and cannot be cancelled by a
regular quadratic form. Thus its coefficient again
vanishes at the origin and the surface equation belongs
to $I^2$. The determinant corrections are four on the
cusp branch and two on the smooth branch $(0,0,t)$.
The total determinant degree is again four.
Dual lifts on the cusp branch are
$t^{-2}\partial_z$ and $\frac12t^{-3}\partial_y$;
on the smooth branch they are
$t^{-1}\partial_x,t^{-1}\partial_y$.
The same product derivative formula bounds the cubic
poles by six and two respectively, of total degree eight.

\noindent\textbf{Step 10.} In either non-lci case let $R\subset N$ be the saturated
radical of the regular rank-one quadratic form and let
$Q=N/R$. Its nonzero square coefficient is a section of
$Q^{-2}\otimes M$, where $M=\OO_Y(S)|_C$ has degree
zero. Hence $\deg Q\le0$ and $\deg R=4-\deg Q\ge4$.
The nonzero intrinsic cubic of a generic $A_2$ is a
section of $R^{-3}\otimes M$ with poles of total degree
at most ten, or eight. It would give respectively
$3\deg R\le10$ or $3\deg R\le8$, both impossible.
This excludes the two non-lci cases, without assuming
that the conductor-image scheme is reduced.

\noindent\textbf{Step 11.} We may consequently use the lci conormal bundle along $H$.
Let $\Xi$ be its normalization conductor divisor, of degree
$2\delta$. The determinant formula and the
conductor-cleared intrinsic cubic from \ref{cdp:st:A.5.5} give
\[
\begin{gathered}
 \det N=K_C(\Xi)\otimes h_C^*K_Y^{-1},\\
 q\in H^0(C,Q^{-2}\otimes M),\qquad
 k\in H^0(C,R^{-3}\otimes M(2\Xi)),\\
 0\ne q^3k^2\in
 H^0\bigl(C,(\det N)^{-6}\otimes M^5(4\Xi)\bigr),\\
 \deg\bigl((\det N)^{-6}\otimes M^5(4\Xi)\bigr)
                         =12-4\delta .
\end{gathered}                                  \tag{C.RA.13}\label{cdp:eq:C.RA.13}
\]
The section is nonzero because the generic transverse
singularity is a cusp.

\noindent\textbf{Step 12.} We give the local comparison that strengthens the zero
count at a smooth point of $H$. A simple zero of $q$ at
which $k$ is nonzero is impossible in a torus smoothing.
Choose coordinates with $H=\{x=y=0\}$, radical direction
$\partial_x$, and coordinate $z$ along $H$.
The nonzero cubic permits Weierstrass preparation in $x$:
\[
 f_S=U\bigl(x^3+a_2(y,z)x^2+a_1(y,z)x+a_0(y,z)\bigr).
 \tag{C.RA.14}\label{cdp:eq:C.RA.14}
\]
The critical-curve equations and the rank-one Hessian
along it give $a_2\in(y)$ and $a_1,a_0\in(y^2)$.
After depressing the cube by $x'=x+a_2/3$, this is
\[
 f_S=U'\bigl(x'^3+y^2(B_1(y,z)x'+B_0(y,z))\bigr).
 \tag{C.RA.15}\label{cdp:eq:C.RA.15}
\]
The simple zero of the Hessian coefficient says
$\partial_zB_0(0,0)\ne0$. Thus
$z'=B_1(y,z)x'+B_0(y,z)$ is an actual holomorphic
coordinate with $x',y$, and the reduced equation is a
unit times $x'^3+y^2z'$. The full fibre equation is a
unit times its $m$th power. \ref{cdp:st:A.5.15} excludes precisely
this Hirzebruch germ, including that unit and multiplicity.

\noindent\textbf{Step 13.} Every point of $C$ above the smooth part of $H$ where a
vertical conductor component meets it has $q=0$.
Indeed a nonzero Hessian would, by \ref{cdp:st:A.5.17}, give a constant
transverse cusp product with smooth normalization and pure
conductor $2C$; its minimal normalization resolution could
have no such vertical component. At a zero of $q$ the
cubic is holomorphic. If it also vanishes, the section
$q^3k^2$ has zero order at least five. If it does not
vanish, \eqref{cdp:eq:C.RA.14}--\eqref{cdp:eq:C.RA.15} force the zero of $q$ to
have order at least two, and the order is at least six.

\noindent\textbf{Step 14.} Each of the $r$ original ruling fibres in \eqref{cdp:eq:C.RA.1} has
a positive-boundary total-transform component meeting $C$,
at a different point of $C$. This remains true at each
blowup: a blowup at its intersection with the coefficient-two
section inserts an exceptional curve of positive coefficient.
Using \eqref{cdp:eq:C.RA.6} and the degree in \eqref{cdp:eq:C.RA.13} now gives
\[
 \begin{array}{c|c|c|c}
 \delta& s\text{ at most}&
 \text{bad points above }H_{\rm reg}\text{ at most}&r\text{ at most}\\ \hline
 0&0&2&2\\
 1&2&1&3\\
 2&3&0&3 .
 \end{array}                                    \tag{C.RA.16}\label{cdp:eq:C.RA.16}
\]
This proves the asserted bound on $r$.

\noindent\textbf{Step 15.} Finally (A.RT.12) and (A.RT.16) give
$e(Z\setminus D_{\rm red})=2-r+N_0$.
The forest argument of (A.RT.21) applies to this entire
star class, whether its minimal complement is negative
or not, and gives $e(G)\ge2$.
Together with $r\le3$, \eqref{cdp:eq:C.RA.2} now gives directly
\[
                 e(S)\ge4-r+N_0\ge1+N_0.
 \tag{C.RA.17}\label{cdp:eq:C.RA.17}
\]
\end{proof}
